%% Produit par outils/generer_interpretation.py — ne pas modifier à la main.
\textbf{1.} A~: \emph{largo}, \emph{andante}, \emph{moderato}, \emph{allegro}, \emph{presto}~; B~: \emph{adagio}, \emph{andante}, \emph{allegretto}, \emph{vivace}, \emph{prestissimo}.

\textbf{2.} \emph{Allegro ma non troppo}~: allègre, mais pas trop~; \emph{Andante con moto}~: en allant, avec mouvement~; \emph{Poco a poco crescendo}~: en croissant peu à peu~; \emph{Più mosso}~: plus animé~; \emph{Dolce e cantabile}~: doux et chantant~; \emph{Sempre staccato}~: toujours détaché~; \emph{Molto ritardando}~: en retardant beaucoup.

\textbf{3.} 1~:~\textbf{\emph{p}}, \emph{piano} (doux)~; 2~:~\emph{staccato} (détaché)~; 3~:~soufflet ouvrant, \emph{crescendo} (en croissant)~; 4~:~\textbf{\emph{f}}, \emph{forte} (fort)~; 5~:~liaison, \emph{legato} (lié)~; 6~:~accent (attaque renforcée)~; 7~:~\emph{marcato} (marqué)~; 8~:~soufflet fermant, \emph{decrescendo} (en décroissant)~; 9~:~\emph{tenuto} (tenu)~; 10~:~\textbf{\emph{pp}}, \emph{pianissimo} (très doux).

\textbf{4.} 1, 2, 1, 2, 3, 4, 5, 6, 3, 4.
